% -*- coding: utf-8 -*-
% Latex代码与渲染效果同时展示.tex
% Latex代码与渲染效果同时展示
\documentclass[UTF8]{ctexart}
\usepackage{geometry}

\geometry{a4paper,scale=0.8}

\usepackage{amsmath}
\usepackage{xcolor}
\usepackage{minted}
\usepackage{tcolorbox}
\tcbuselibrary{listings,skins}

% 定义 latexexample 环境
\newtcblisting{latexexample}[1][]{
    colback=gray!10,
    colframe=gray!50,
    listing side text,  % 左侧代码，右侧效果
    title={LaTeX 示例},
    fonttitle=\bfseries,
    #1
}

\title{\LaTeX 代码与渲染效果同时展示}
\author{埃及猪肉}
\date{\today}

\begin{document}

\maketitle
\tableofcontents

\section{LaTeX 示例}

\begin{latexexample}
Good: One, two, three\ldots

Bad: One, two, three...
\end{latexexample}

\section{实现细节}

\subsection{实现方法}

\subsubsection{完整代码示例}

以下是完整的实现代码，包含导言区配置和正文使用：

\begin{tcolorbox}[colback=gray!5,colframe=gray!50,title={完整代码}]
\begin{verbatim}
% 导言区
\documentclass[UTF8]{ctexart}
\usepackage{tcolorbox}
\tcbuselibrary{listings,skins}

\newtcblisting{latexexample}[1][]{
    colback=gray!10,
    colframe=gray!50,
    listing side text,
    title={LaTeX 示例},
    fonttitle=\bfseries,
    #1
}

% 正文区
\begin{document}
\begin{latexexample}
Good: One, two, three\ldots

Bad: One, two, three...
\end{latexexample}
\end{document}
\end{verbatim}
\end{tcolorbox}

\subsubsection{分步说明}

在导言区定义环境，然后在正文中使用：

\begin{enumerate}
    \item \textbf{导言区定义}
    \begin{tcolorbox}[colback=gray!5,colframe=gray!50,title={导言区代码}]
\begin{verbatim}
\newtcblisting{latexexample}[1][]{
    colback=gray!10,
    colframe=gray!50,
    listing side text,
    title={LaTeX 示例},
    fonttitle=\bfseries,
    #1
}
\end{verbatim}
    \end{tcolorbox}

    \item \textbf{正文中使用}
    \begin{tcolorbox}[colback=gray!5,colframe=gray!50,title={正文代码}]
\begin{verbatim}
\begin{latexexample}
Good: One, two, three\ldots

Bad: One, two, three...
\end{latexexample}
\end{verbatim}
    \end{tcolorbox}

    \item \textbf{渲染效果}

    \begin{latexexample}
Good: One, two, three\ldots

Bad: One, two, three...
    \end{latexexample}
\end{enumerate}

\subsection{核心原理}

本功能依赖于 \texttt{tcolorbox} 宏包的 \texttt{listings} 库，通过 \texttt{newtcblisting} 命令创建自定义环境。

\texttt{listing side text} 选项将内容区域分为两部分：
\begin{itemize}
    \item 左侧：显示源代码（由 listings 处理）
    \item 右侧：显示渲染效果（由 LaTeX 编译）
\end{itemize}

\subsection{环境定义解析}

\begin{itemize}
    \item \texttt{colback=gray!10}：设置背景颜色为浅灰色
    \item \texttt{colframe=gray!50}：设置边框颜色为中灰色
    \item \texttt{listing side text}：启用左右分栏布局
    \item \texttt{title=\{...\}}：设置标题
    \item \texttt{fonttitle=\textbackslash bfseries}：标题加粗
    \item \texttt{\#1}：允许传递额外参数进行自定义
\end{itemize}

\subsection{使用限制}

\begin{itemize}
    \item 不支持包含 \texttt{\$...\$} 的复杂数学公式
    \item 不支持 \texttt{\textbackslash begin\{...\}} 环境嵌套
    \item 不支持导言区命令（如 \texttt{\textbackslash usepackage}）
    \item 适合展示简单的文本和格式示例
\end{itemize}


\end{document}
